% META: {"id":"1SPE-PRODSCAL-EX-033","chapitre":"1SPE-PRODUIT-SCALAIRE","type_objet":"exercice","capacites":["1SPE-PRODUIT-SCALAIRE-C4"],"capacites_codes":["C4"],"methodes":["M4"],"parcours":1,"competences":["calculer"],"duree_min":5,"mode_creation":"ex_nihilo","sources_inspiration":[],"fichier_tex":"chapitres/1SPE-PRODUIT-SCALAIRE/exercices/1SPE-PRODSCAL-EX-033.tex","corrige_tex":"chapitres/1SPE-PRODUIT-SCALAIRE/corriges/1SPE-PRODSCAL-CO-033.tex","status":"approved","origin":"generated","status_history":["generated","audited","accepted","approved"],"release_acceptance":"RELEASE_OWNER_BATCH_ACCEPTANCE","acceptance_closure_digest":"sha256:9b3ccf9a81c5520fbb7e03b7b2d3e7bf2057a02834b6d908bf3be5f49f3b553f","acceptance_receipt_digest":"sha256:4fad86f0282e0fa4e0419cca1ad6379ae7af5a280c04a4a90880f90a1c044242"}
% BEGIN-VERIFY
% from sympy import *
% # A=(0,0), B=(4,0), C=(1,3): aire = 1/2 * |4*3 - 0*1| = 6
% assert abs(4*3 - 0*1) // 2 == 6
% END-VERIFY
\begin{exercice}{1SPE-PRODSCAL-EX-033}{1}{6}
Calculer l'aire du triangle $ABC$ avec $A(0; 0)$, $B(4; 0)$ et $C(1; 3)$.
\end{exercice}
